\begin{lemma}
\label{lemma-ML-countable-colimit}
Let $M$ be an $R$-module.  Write $M = \colim_{i \in I} M_i$ where $(M_i,
f_{ij})$ is a directed system of finitely presented $R$-modules.  If $M$ is
Mittag-Leffler and countably generated, then there is a directed countable
subset $I' \subset I$ such that $M \cong \colim_{i \in I'} M_i$.
\end{lemma}

\begin{proof}
Let $x_1, x_2, \ldots$ be a countable set of generators for $M$.  For each $x_n$
choose $i \in I$ such that $x_n$ is in the image of the canonical map $f_i:
M_i \to M$; let $I'_{0} \subset I$ be the set of all these $i$.  Now
since $M$ is Mittag-Leffler, for each $i \in I'_{0}$ we can choose $j \in I$
such that $j \geq i$ and $f_{ij}: M_i \to M_j$ factors through
$f_{ik}: M_i \to M_k$ for all $k \geq i$  (condition (3) of Proposition
\ref{proposition-ML-characterization}); let $I'_1$ be the union of $I'_0$ with
all of these $j$.  Since $I'_1$ is a countable, we can enlarge it to a
countable directed set $I'_{2} \subset I$.  Now we can apply the same procedure
to $I'_{2}$ as we did to $I'_{0}$ to get a new countable set $I'_{3} \subset
I$.  Then we enlarge $I'_{3}$ to a countable directed set $I'_{4}$.  Continuing
in this way---adding in a $j$ as in Proposition
\ref{proposition-ML-characterization} (3) for each $ i \in I'_{\ell}$ if $\ell$
is odd and enlarging $I'_{\ell}$ to a directed set if $\ell$ is even---we get a
sequence of subsets $I'_{\ell} \subset I$ for $\ell \geq 0$.  The union $I' =
\bigcup I'_{\ell}$ satisfies:
\begin{enumerate}
\item $I'$ is countable and directed;
\item each $x_n$ is in the image of $f_i: M_i \to M$ for some $i
\in I'$;
\item if $i \in I'$, then there is $j \in I'$ such that $j \geq i$ and $f_{ij}:
M_i \to M_j$ factors through $f_{ik}: M_i \to M_k$ for all
$k \in I$ with $k \geq i$.  In particular $\Ker(f_{ik}) \subset \Ker(f_{ij})$
for $k \geq i$.
\end{enumerate}
We claim that the canonical map $\colim_{i \in I'} M_i \to
\colim_{i
\in I} M_i = M$ is an isomorphism.  By (2) it is surjective.  For injectivity,
suppose $x \in \colim_{i \in I'} M_i$ maps to $0$ in $\colim_{i \in
I} M_i$.
Representing $x$ by an element $\tilde{x} \in M_i$ for some $i \in I'$, this
means that $f_{ik}(\tilde{x}) = 0$ for some $k \in I, k \geq i$.  But then by
(3) there is $j \in I', j \geq i,$ such that $f_{ij}(\tilde{x}) = 0$.  Hence $x
= 0$ in $\colim_{i \in I'} M_i$.
\end{proof}
